%% start of file `template.tex'.
%% Copyright 2006-2013 Xavier Danaux (xdanaux@gmail.com).
%
% This work may be distributed and/or modified under the
% conditions of the LaTeX Project Public License version 1.3c,
% available at http://www.latex-project.org/lppl/.


\documentclass[10pt,letterpaper,sans]{moderncv}
\moderncvstyle[nosymbols]{banking}
\moderncvcolor{blue}
\usepackage{lmodern}
\usepackage{hyperref}
\hypersetup{colorlinks=true, urlcolor=color1, linkcolor=color1}
\usepackage{fancyhdr}
\pagestyle{fancy}
\usepackage[utf8]{inputenc}
\usepackage[margin=0.8in, top=1in]{geometry}
\usepackage{amsmath}
\providecommand{\sl}{\slshape}
\AtBeginDocument{\microtypesetup{expansion=false}}
\colorlet{firstnamecolor}{color1}
\colorlet{lastnamecolor}{color1}
\colorlet{namecolor}{color1}
\colorlet{addresscolor}{color1}
\renewcommand*{\emailsymbol}{}
\renewcommand*{\homepagesymbol}{}
\renewcommand*{\sectionfont}{\sffamily\Large\bfseries\upshape}
\makeatletter
\renewcommand*{\makecvhead}{%
  \recomputecvlengths
  \makehead
  \vspace{1em}%
  \par}
\AtBeginDocument{%
  \renewcommand{\rmdefault}{\sfdefault}%
  \renewcommand{\familydefault}{\sfdefault}%
  \normalfont\sffamily
  \colorlet{firstnamecolor}{color1}%
  \colorlet{lastnamecolor}{color1}}
\makeatother
% personal data
\name{Mithilesh}{Vaidya}
% \title{Curriculum Vitae}
% \title{CV}                              % optional, remove / comment the line if not wanted
%\address{my address, line 1, line 2, line 3, postcode}{}{}% optional, remove / comment the line if not wanted; the "postcode city" and and "country" arguments can be omitted or provided empty
%\phone[mobile]{+44 12345 12345}                   % optional, remove / comment the line if not wanted
%\phone[fixed]{01234 123456}                    % optional, remove / comment the line if not wanted
%\phone[fax]{+3~(456)~789~012}                      % optional, remove / comment the line if not wanted
  
\email{vaidya.mithilesh@gmail.com}                               % optional, remove / comment the line if not wanted
\homepage{methi1999.github.io}                       % optional, remove / comment the line if not wanted
% \phone[Phone]{+1-425-900-7799}
% \extrainfo{+1 425-900-7799}
\extrainfo{\href{https://www.linkedin.com/in/mithileshvaidya/}{LinkedIn}}
% \extrainfo{\textbf{Final year student at IIT Bombay | Applying for a Master's in Computer Science}}
% \extrainfo{\textbf{Final year student at IIT Bombay | User ID: mithilesh.vaidya@iitb.ac.in}}
% \extrainfo{\textbf{Final year student at IIT Bombay}}
%----------------------------------------------------------------------------------
%            content
%----------------------------------------------------------------------------------

\begin{document}
\renewcommand{\labelitemi}{$\bullet$}
\makecvtitle
%%%%%%%%%%%%%%%%%
%%%%%%%%%%%%%%%%%
% \section{\textsc{Research Interests}}
% Machine Learning, Speech Processing, Prosody
%%%%%%%%%%%%%%%%%
%%%%%%%%%%%%%%%%%
\section{Education}

\begin{itemize}
\item \textbf{Georgia Institute of Technology}, Atlanta, USA \hfill{[2022 - 2024]}
\begin{itemize}
    \item Master's in Computer Science (Specialisation: Machine Learning)
\end{itemize}
\vspace{10pt}
\item \textbf{Indian Institute of Technology Bombay}, Mumbai, India \hfill{[2017 - 2022]}
\begin{itemize}
    \item Bachelor and Master of Technology in Electrical Engineering, \textbf{CPI 9.51/10}
    \item \textbf{Minor} in Computer Science and Engineering
\end{itemize}
\end{itemize}
\vspace{-10pt}
%%%%%%%%%%%%%%%%%
%%%%%%%%%%%%%%%%%
\section{Honors and Awards}
\begin{itemize}
\item \textbf{Institute Silver Medal} at IIT Bombay for best academic standing in department \hfill{\small [2022]}
\item \textbf{Excellence in Academics} Award, for excellent academic performance (Top 3/121) \hfill{\small [2022]}
\item \textbf{Undergraduate Research Award} (URA03) for outstanding thesis contributions \hfill{\small [2022]}
\item Bhavesh Gandhi Memorial Prize for standing 1st in the CSP specialization
\item \textbf{JN Tata} Scholarship for pursuing higher education abroad\hfill {\small[2022]}
\item \textbf{AP grade} in \textbf{Control Systems} course for \textbf{exceptional performance} \hfill{\small [2020]}
\item \textbf{All India Rank of 388} in JEE Advanced 2017 among 2,00,000 candidates \hfill {\small [2017]}
% \item Awarded Certificate of Merit for being among the state \textbf{top 1\%} in \textbf{NSEA} and \textbf{NSEP}, \\organised by India Association of Physics Teachers \hfill{\small [2016]}
\item \textbf{Kishore Vaigyanik Protsahan Yojana (KVPY)} Fellowship by Govt. of India \hfill{\small [2015]}
\item \textbf{National Talent Search Examination (NTSE)} scholarship by Govt. of India\hfill{\small [2015]}
\item \textbf{Silver medal} in Homi Bhabha Young Scientist Examination \hfill{\small [2011]}
\end{itemize}
\vspace{-10pt}
%%%%%%%%%%%%%%%%%
%%%%%%%%%%%%%%%%%
\section{Publications and Pre-prints}
\begin{itemize}
    \item A. Luebs, \textbf{Mithilesh Vaidya}, I. Kumar, S. Badam, S. W. Bailey, M. Bendel, J. Sotelo and X. He. ``PoDAR: Power-Disentangled Audio Representation for Generative Modeling.'' arXiv preprint \href{https://arxiv.org/abs/2605.10084}{arXiv:2605.10084} (2026).
    \item \textbf{M. Vaidya}, C. Rodgers and A. Wu. ``\href{https://methi1999.github.io/assets/pdf/switching_subspace_model_neurips.pdf}{Switching Subspace Model for Neural Population Analysis}'' (2024).
    \item \textbf{Vaidya, Mithilesh}, Binaya Kumar Sahoo, and Preeti Rao. ``Deep Learning for Assessment of Oral Reading Fluency.'' arXiv preprint \href{https://arxiv.org/abs/2405.19426}{arXiv:2405.19426} (2024).
    \item \textbf{M. Vaidya}, K. Sabu and P. Rao. ``Deep Learning for Prominence Detection In Children's Read Speech.'' arXiv preprint \href{https://arxiv.org/abs/2110.14273}{arXiv:2110.14273} (2021).
    \item Sabu, Kamini, \textbf{Mithilesh Vaidya}, and Preeti Rao. ``CNN Encoding of Acoustic Parameters for Prominence Detection.'' arXiv preprint \href{https://arxiv.org/abs/2104.05488}{arXiv:2104.05488} (2021).
\end{itemize}
\vspace{-10pt}
%%%%%%%%%%%%%%%
%%%%%%%%%%%%%%%
\section{Professional Experience}
\begin{itemize}
\item \textbf{AI Researcher} \hfill{\sl June'24 - Present}\\
\textit{\href{https://www.descript.com/research}{Descript}, San Francisco} \hfill {\sl Research Engineer}\\
\vspace{-12pt}
\begin{itemize}
\item Co-led development of \href{https://descriptinc.github.io/uvm-v2/}{\textbf{Audio Regenerate}}, a production latent-inpainting system for seamless word-level audio editing in Descript
\item Trained both stages of the Audio Regenerate pipeline: a continuous neural audio codec and a flow-matching generator for masked latent inpainting
\item Developed \href{https://descriptinc.github.io/video-regenerate/}{\textbf{Video Regenerate}}, an audio-driven lip-sync model for video editing and translation, including a contrastive model for audio-video alignment
\item Co-authored \href{https://arxiv.org/abs/2605.10084}{\textbf{PoDAR}}, a power-disentangled audio representation that improves latent modelability and accelerates generator convergence by $\mathbf{2\times}$
\item Improved the data preprocessing pipeline for training generative video models on large-scale datasets
\end{itemize}
\end{itemize}
\begin{itemize}
\item \textbf{LLM-powered Natural Language Interface} \hfill{\sl May'23 - Aug'23}\\
\textit{Qualcomm, San Diego | Guide: Mr. Vasudev Nayak, Principal Engineer} \hfill {\sl Machine Learning Internship}\\
\vspace{-12pt}
\begin{itemize}
\item Designed an end-to-end hierarchical modular pipeline for structured parsing using LLMs such as MPT-7B
\item Developed techniques for reducing hallucinations and compared performance across various open-source LLMs
\item Evaluated performance of NVIDIA Nemo and OpenAI Whisper for STT and machine translation
\end{itemize}
\vspace{5pt}
\item \textbf{Verification of FPGA-based High Frequency Trading Platform} \hfill{\sl Apr'20 - June'20}\\
\textit{APT Portfolio Pvt. Ltd. | Guide: Mr. Vivek Pannikar, Senior Verification Engineer} \hfill {\sl Internship}\\
\vspace{-10pt}
\begin{itemize}
\item Implemented \textbf{DPI}, a protocol for exchanging data between SystemVerilog and other languages, for \textbf{speeding up verification} of testbenches using Cocotb, Quartus and Riviera by \textbf{3x}
\item Used Python \textbf{metaclasses} for automatically generating Python, SystemVerilog and C DPI header and implementation files from high-level JSON inputs
% \item \textbf{Speeded} up verification of certain transaction types by \textbf{upto 3x}
\end{itemize}
\vspace{5pt}
%%%%%%%%%%%%%%%
\item \textbf{Autonise AI} \hfill{\sl Sep'18 - May'19}\\
\textit{Machine Learning Startup} \hfill {\sl Co-Founder}\\
\vspace{-10pt}
\begin{itemize}
\item Implemented \textbf{PixelLink} and a GRU for word-level text detection, \textbf{invariant} to font size, colour, background, orientation, etc. and demonstrated an accuracy of \textbf{74\%} on a proprietary dataset of documents like Aadhar Card, Passport, Driving Licence, etc.
\item Implemented a robust model with a \textbf{UNet} backbone for \textbf{segmenting} out spots, patches and wrinkles in selfies and exposed it through \textbf{AWS for demonstration}
\end{itemize}
\end{itemize}
%%%%%%%%%%%%%%%%%
%%%%%%%%%%%%%%%%%
\section{Research Experience}
\begin{itemize}
%%%%%%%%%%%%%%%
\item \textbf{\href{https://methi1999.github.io/assets/pdf/switching_subspace_model_neurips.pdf}{Switching Subspace Model for Neural Population Analysis}} \hfill{\sl Sep'22 - May'24}\\
\textit{Guide: Prof. Anqi Wu, GeorgiaTech} \hfill{\sl Research project}\\
\vspace{-10pt}
\begin{itemize}
\item Developed a sequential VAE model to capture multiple latent dynamics throughout a trial
\item Introduced unimodal Gaussian Process prior to improve interpretability
\item Introduce supervision using CNNs in order to preserve behavioral variable of interest
\end{itemize}
\vspace{5pt}
\item \textbf{Assessing Comprehensibility of Children's Read Speech} \hfill{\sl  May'21 - June'22}\\
\textit{Guide: Prof. Preeti Rao, IIT Bombay | \href{https://arxiv.org/abs/2405.19426}{arXiv:2405.19426}} \hfill {\sl Master's Thesis II}\\
\vspace{-10pt}
\begin{itemize}
% \item Extracted acoustic features from \textbf{raw waveforms} using various deep convolutional architectures
% \item Experimenting with \textbf{transfer learning} by using pre-trained models trained on Emotion Recognition datasets
% \item Exploring Attention mechanisms for \textbf{fusing} lexical and acoustic features
% \item Exploring self-supervised learning for extracting \textbf{prosody embeddings} from a vast unlabelled dataset
\item Implemented a Wav2vec2.0-based end-to-end model for oral reading fluency assessment
\item Outperformed RFC baseline operating on hand-crafted features by 0.06 (absolute Pearson)
\item Probed the internal representations for presence of knowledge-based features using MLP probes
% \item Concatenated hand-crafted features at various hierarchies to further improve performance
\item Implemented a multi-task learning framework for exploiting hand-crafted features during training
\item Proposed a self-supervised learning framework for utilising an unlabelled dataset
\end{itemize}
\vspace{5pt}
%%%%%%%%%%%%%%%%%
\item \textbf{Prominence Detection in Children's Read Speech} \hfill{\sl  Jan'21 - Oct'21}\\
\textit{Guide: Prof. Preeti Rao, IIT Bombay | \href{https://arxiv.org/abs/2110.14273}{arXiv:2110.14273}} \hfill {\sl Master's Thesis I}\\
\vspace{-10pt}
\begin{itemize}
\item Replaced a Random Forest Classifier baseline with a \textbf{CRNN} framework for predicting the degree of prominence for each word in children’s read speech
\item Explored inputs at various hierarchies: raw waveforms, acoustic contours and word-level aggregates
\item Demonstrated an improvement in the acoustic features extracted from word segments using Sinc convolution
\item Exploited phrase boundary labels in various \textbf{multi-task learning} paradigms
\item Used part-of-speech tags and various \textbf{NLP} embeddings such as \textbf{GloVe} and \textbf{BERT} for incorporating complementary lexical information
\end{itemize}
\vspace{5pt}
%%%%%%%%%%%%%%%
\item \textbf{Character Animation from Video for Blender} \hfill{\sl  July'21 - Dec'21}\\
\textit{Guide: Prof. Parag Chaudhuri, IIT Bombay} \hfill{\sl Research Project}\\
\vspace{-10pt}
\begin{itemize}
\item Working on a Blender plugin consisting of an \textbf{integrated pipeline} for extracting \textbf{3D human pose} from a video using various deep learning backends and \textbf{retargeting} it to a rigged character in Blender
\item Added VIBE and MediaPipe to the pose extraction backend of the plugin
\item Explored a \textbf{self-supervised} graph neural network framework for dynamic mapping of animations from source to dissimilar target skeletons
\end{itemize}
\vspace{5pt}
%%%%%%%%%%%%%%%
\item \textbf{SIRD Dynamics} \hfill{\sl Aug'20 - Dec'20}\\
\textit{Guide: Prof. Sharayu Moharir, IIT Bombay} \hfill{\sl Research Project}\\
\vspace{-10pt}
\begin{itemize}
\item Studied the \textbf{SIRD} model which is widely used for studying the outbreak of epidemics
\item Simulated the model with various underlying \textbf{network topologies} in place of random mixing
\item Formulated multiple mathematical models for calculation of \textbf{precise dynamics}
\end{itemize}
%%%%%%%%%%%%%%%
% \item \textbf{Keyword Spotting} \hfill{\sl  June'19 - June'20}\\
% \textit{Guide: Prof. Preeti Rao, IIT Bombay} \hfill{\sl Research Project}\\
% \vspace{-10pt}
% \begin{itemize}
% \item Trained various deep architectures on MFCC features for spotting a keyword in an audio file
% \item Experimented with triplet loss for low-resource setting
% \item hoishdohs
% \end{itemize}

\end{itemize}
\vspace{-10pt}
%%%%%%%%%%%%%%%
%%%%%%%%%%%%%%%
\section{Key Projects}
\begin{itemize}
\item \textbf{WhichLORA} | Course Project: Systems for AI: LLMs \hfill{\sl  Jan'24 - Apr'24}\\
\vspace{-12pt}
\begin{itemize}
\item Developed an end-to-end two-stage framework for automatic generation of LoRA variants \\ based on user constraints for fine-tuning and inference compute and latency
\item Profiled various hardware configurations such as NVIDIA A100, H100 and RTX Quadro 6000
\item Developed a memory predictor for predicting HBM usage during inference given a LLM configuration
\end{itemize}
\vspace{5pt}

\item \textbf{Causal Relationships in Diabetes Prediction} | Course Project: Machine Learning \hfill{\sl  Aug'23 - Dec'23}\\
\vspace{-12pt}
\begin{itemize}
\item Conducted EDA of the 21 features using violin plots, heatmaps and bar charts
\item Implemented various baselines such as SVM, Random Forests and XGBoost for predicting diabetes
\item Uncovered various causal relationships by designing various DAGs
\end{itemize}
\vspace{5pt}

\item \textbf{Compiler Optimizations for speeding up Capsules} | Course Project: Parallelizing Compilers \hfill{\sl  Mar'23 - May'23}\\
\vspace{-12pt}
\begin{itemize}
\item Implemented compiler optimization techniques such as unroll and jam, scalar replacement and\\ loop reordering in CUDA for reducing execution time of the core capsule operation
\item Explained the \textbf{4x speedup} in CUDA execution time by examining the profiler output
\end{itemize}
\vspace{5pt}
\item \textbf{Legendre Memory Unit} | Course Project: Advanced Machine Learning \hfill{\sl Jan'21 - May'21}\\
\vspace{-10pt}
\begin{itemize}
\item Implemented and analysed the performance of Legendre Memory Units (LMU), \textbf{an improved sequential model}, on various tasks and datasets such as JSB Chorales, Mackey-Glass dynamics, etc.
\item Suggested modifications to the \textbf{core equations} by studying various \textbf{basis} functions
\end{itemize}
\vspace{5pt}
%%%%%%%%%%%%%%%
\item \textbf{Audio Steganography} | Course Project: Automatic Speech Recognition \hfill{\sl Jan'21 - May'21}\\
\vspace{-10pt}
\begin{itemize}
\item Exploited \textbf{adversarial attacks} on ASR systems for hiding any given sequence of tokens in any audio file
\item Analysed performance as a function of various token sequence properties such as \textbf{length and perplexity}
\item Demonstrated \textbf{high PESQ scores} which indicate low perceptibility of deviation from original audio
\end{itemize}
\vspace{5pt}
%%%%%%%%%%%%%%%
\item \textbf{Video Toonification} | Course Project: Digital Image Processing \hfill{\sl Aug'20 - Dec'20}\\
\vspace{-10pt}
\begin{itemize}
\item Used \textbf{Mean Shift Segmentation} across both time and spatial dimensions for toonification of videos
\item Benchmarked results with standard techniques such as Bilateral Filtering
\end{itemize}
\vspace{5pt}
%%%%%%%%%%%%%%%
\item \textbf{Auction Theory} | Course Project: Game Theory \hfill{\sl Aug'20 - Dec'20}\\
\vspace{-10pt}
\begin{itemize}
\item Studied various models in Auction Theory such as \textbf{Vickrey Auction} and \textbf{First Price Sealed Bid} Auctions
\item Discussed equilibrium and optimal auction design analysis
\end{itemize}
\vspace{5pt}
%%%%%%%%%%%%%%%
\item \textbf{FMX Rendering and Animation} | Course Project: Computer Graphics \hfill{\sl Nov'20 - Dec'20}\\
\vspace{-10pt}
\begin{itemize}
\item Designed and rendered an FMX track with obstacles of varying shapes such as \textbf{cylinders and ramps}
\item Designed, rendered and animated a rider and a motorbike on the track using \textbf{keyframing}
\item Employed \textbf{Phong Shading, Texture Mapping} and used a \textbf{Skybox} for a realistic look
\end{itemize}
\vspace{5pt}
%%%%%%%%%%%%%%%
\item \textbf{Pipelined RISC Processor} | Course Project: Microprocessors \hfill{\sl Oct'19 - Nov'19}\\
\vspace{-10pt}
\begin{itemize}
\item Designed a 16-bit, 8-register, 6-stage \textbf{Pipelined} RISC processor in VHDL
% \item Tested it on an FPGA using a custom testbench covering all corner cases
\item Employed \textbf{Branch Prediction} and \textbf{Hazard Mitigation} techniques for optimizing the performance
\end{itemize}
\vspace{5pt}
%%%%%%%%%%%%%%%
\item \textbf{FindIt} | Self Project: Audio Fingerprinting \hfill{\sl May'19 - June'19}\\
\vspace{-10pt}
\begin{itemize}
\item FindIt is a Python program for identifying a song given a \textbf{short noisy segment}, similar to Shazam
\item An \textbf{audio fingerprint} consisting of constellations of major time-frequency peaks is stored in a hash table
% \item The music discovery program has \textbf{30+ stars} on GitHub
\end{itemize}
\vspace{5pt}
%%%%%%%%%%%%%%%
\item \textbf{Handwriting Recognition Pen} | Summer Project \hfill{\sl May'18 - June'18}\\
\vspace{-10pt}
\begin{itemize}
\item Built a pen which can instantly convert \textbf{handwriting strokes} on ordinary paper into text
\item Designed the pen from scratch in AutoCAD and \textbf{3D printed} it
\item Generated own training data for each letter using a custom OpenCV script
\end{itemize}
\vspace{5pt}
%%%%%%%%%%%%%%%
\item \textbf{Furniture Classification} | Kaggle Competition \hfill{\sl May'18 - June'18}\\
\vspace{-10pt}
\begin{itemize}
\item Participated in the iMaterialist Challenge organised by Malong Technologies and CVPR 2018
\item Implemented \textbf{ResNet in PyTorch} for classification of furniture images into \textbf{128 classes}, each class \\ containing around 1500 training images with low inter-class variation
        % \item Fine-tuned the hyperparameters of the model for improved accuracy under different lighting, angles,\\backgrounds, and levels of occlusion
\item Achieved an accuracy of 87.4\% and ranked \textbf{30} among 428 teams across the globe
\end{itemize}

\end{itemize}
\vspace{-10pt}
%%%%%%%%%%%%%%%
%%%%%%%%%%%%%%%
\section{Technical Skills}
\vspace{3pt}
\begin{itemize}
\item \textbf{Programming Languages:} Python, C++, C, Bash, Verilog, VHDL, OpenGL, SQLite
\vspace{3pt}
\item \textbf{Softwares:} Matlab, Arduino, \LaTeX, Blender, VHDL, AutoCad, Solidworks, Android Studio
\vspace{3pt}
\item \textbf{Data Science:} PyTorch, Pandas, Numpy, OpenCV, TensorFlow, MATLAB
\end{itemize}
\vspace{-10pt}
%%%%%%%%%%%%%%%
%%%%%%%%%%%%%%%
\section{Key Coursework}
\vspace{3pt}
\begin{itemize}
\item \textbf{Computer Science:} Logic, Compilers, Systems for AI, Game AI, Automatic Speech Recognition, Advanced Machine Learning, Computer Graphics, Foundations of Intelligent and Learning Agents, Data Structures and Algorithms, Digital Image Processing, Network Security, Computer Networks, Operating Systems, Data and Visual Analysis
\vspace{3pt}
\item{\textbf{Electrical Engineering:} Digital Communication, Digital Systems, Digital Signal Processing, Data Analysis and Interpretation, Control Systems, Information Theory and Coding, Markov Chains, Microprocessors, Microelectronics, Foundations of VLSI CAD, Power Systems}
\vspace{3pt}
\item{\textbf{Miscellaneous:} Non-linear Optimization, Dynamic Algebraic Algorithms, Calculus, Complex Analysis, Linear Algebra, Differential Equations, Biology, Chemistry, Economics, Psychology, Engineering Drawing, Environmental Studies}
\end{itemize}
% \hfill{\sl * To be completed in December 2020}
\vspace{-10pt}
%%%%%%%%%%%%%%%
%%%%%%%%%%%%%%%
\section{Positions of Responsibility}
\begin{itemize}
\item \textbf{Editorial Board Member, Insight} \hfill{\sl Apr'21 - Mar'22}\\
\textit{Insight is IIT Bombay's student media body with over 10,000+ readers}
\begin{itemize}
\item Surveyed the \textbf{effectiveness} of the Faculty Advisor program by taking inputs from both students and faculty and \textbf{suggested various reforms}
\item Initiated a series on \textbf{startups from research labs} at IIT Bombay as part of the \textbf{LinkedIn} team
\item Interviewed authorities and current international students for understanding the causes behind \textbf{poor international representation} at IITB and suggested \textbf{remedies} for the same
\item Awarded \textbf{Special Mention for Journalism} by IIT Bombay for work carried out as part of Insight
\end{itemize}
\vspace{5pt}
%%%%%%%%%%%%%%%
\item \textbf{Department Academic Mentor} \hfill{\sl Apr'21 - June'22}\\
\vspace{-10pt}
\begin{itemize}
\item Selected as part of a 35-member team on the basis of \textbf{ethics, peer-reviews} and an \textbf{interview}
\item \textbf{Mentoring} 6 sophomores in academic and co-curricular activities
\end{itemize}
\vspace{5pt}
%%%%%%%%%%%%%%%
\item \textbf{Teaching Assistant} \hfill{\sl July'21 - Apr'22}\\
\vspace{-10pt}

\begin{itemize}
\item TA for EE679, a graduate-level course on Speech Processing which covers speech production, analysis techniques and applications such as ASR, speech synthesis, etc. (Autumn 2021-22)
\item TA for EE352, a lab course on Digital Signal Processing which covers practical aspects of signal processing algorithms such as FFT on digital signal processors (Spring 2021-22)
% \item In charge of assisting the instructor in conducting the evaluation of the course
\end{itemize}
\end{itemize}
\vspace{-10pt}
%%%%%%%%%%%%%%%
%%%%%%%%%%%%%%%
\section{Extra-Curricular Activities}
\vspace{3pt}
\begin{itemize}
\item National-level quarter-finalist at Bournvita Quiz Contest; appeared on \textbf{National TV} for the same
\item \textbf{Won 2nd prize} in \textbf{Android app development} competition organised by Web and Coding Club
\item Successfully completed a 12-month \textbf{Lawn Tennis} course under National Sports Organisation and represented Hostel 4 in inter-hostel tournaments
\item Attended the \textbf{Vijyoshi} camp conducted by \textbf{IISC}, Bangalore which served as a platform for interaction between bright young students and leading researchers in the field of science
\item Awarded \textbf{Best Outgoing Student} of the year 2014-15 by Nirmala Convent High School
\end{itemize}
\end{document}



% %%%%%%%%%%%%%%%
% %%%%%%%%%%%%%%%
% \section{Key Projects}
% \begin{itemize}
% \item \textbf{Legendre Memory Unit} \hfill{\sl Jan'21 - May'21}\\
% \textit{Instructor: Prof. Sunita Sarawagi} \hfill{\sl Course Project: Advanced Machine Learning}\\
% \vspace{-10pt}
% \begin{itemize}
% \item Implemented and analysed the performance of Legendre Memory Units (LMU), \textbf{an improved sequential model}, on various tasks and datasets such as JSB Chorales, Mackey-Glass dynamics, etc.
% \item Suggested modifications to the \textbf{core equations} by studying various \textbf{basis} functions
% \end{itemize}
% \vspace{5pt}
% %%%%%%%%%%%%%%%
% \item \textbf{Audio Steganography} \hfill{\sl Jan'21 - May'21}\\
% \textit{Instructor: Prof. Preethi Jyothi} \hfill{\sl Course Project: Automatic Speech Recognition}\\
% \vspace{-10pt}
% \begin{itemize}
% \item Exploited \textbf{adversarial attacks} on ASR systems for hiding any given sequence of tokens in any audio file
% \item Analysed performance as a function of various token sequence properties such as \textbf{length and perplexity}
% \item Demonstrated \textbf{high PESQ scores} which indicate low perceptibility of deviation from original audio
% \end{itemize}
% \vspace{5pt}
% %%%%%%%%%%%%%%%
% \item \textbf{Video Toonification} \hfill{\sl Aug'20 - Dec'20}\\
% \textit{Instructor: Prof. Ajit Rajwade \& Prof. Suyash Awate} \hfill{\sl Course Project: Digital Image Processing}\\
% \vspace{-10pt}
% \begin{itemize}
% \item Used \textbf{Mean Shift Segmentation} across both time and spatial dimensions for toonification of videos
% \item Benchmarked results with standard techniques such as Bilateral Filtering
% \end{itemize}
% \vspace{5pt}
% %%%%%%%%%%%%%%%
% \item \textbf{Auction Theory} \hfill{\sl Aug'20 - Dec'20}\\
% \textit{Instructor: Prof. Ankur Kulkarni} \hfill{\sl Course Project: Game Theory}\\
% \vspace{-10pt}

% \begin{itemize}
% \item Studied various models in Auction Theory such as \textbf{Vickrey Auction} and \textbf{First Price Sealed Bid} Auctions
% % \item Discussed \textbf{equilibrium and optimal auction design} analysis
% \end{itemize}
% \vspace{5pt}
% %%%%%%%%%%%%%%%
% \item \textbf{FMX Rendering and Animation} \hfill{\sl Nov'20 - Dec'20}\\
% \textit{Instructor: Prof. Parag Chaudhari} \hfill{\sl Course Project: Computer Graphics}\\
% \vspace{-10pt}

% \begin{itemize}
% \item Designed and rendered an FMX track with obstacles of varying shapes such as \textbf{cylinders and ramps}
% \item Designed, rendered and animated a rider and a motorbike on the track using \textbf{keyframing}
% \item Employed \textbf{Phong Shading, Texture Mapping} and used a \textbf{Skybox} for a realistic look
% \end{itemize}
% \vspace{5pt}
% %%%%%%%%%%%%%%%
% \item \textbf{Pipelined RISC Processor} \hfill{\sl Oct'19 - Nov'19}\\
% \textit{Instructor: Prof. Virendra Singh} \hfill{\sl Course Project: Microprocessors}\\
% \vspace{-10pt}

% \begin{itemize}
% \item Designed a 16-bit, 8-register, 6-stage \textbf{Pipelined} RISC processor in VHDL
% \item Tested it on an FPGA using a custom testbench covering all corner cases
% \item Employed \textbf{Branch Prediction} and \textbf{Hazard Mitigation} techniques for optimizing the performance
% \end{itemize}
% \vspace{5pt}
% %%%%%%%%%%%%%%%
% \item \textbf{FindIt} \hfill{\sl May'19 - June'19}\\
% \textit{Audio Fingerprinting} \hfill{\sl Self Project}\\
% \vspace{-10pt}

% \begin{itemize}
% \item FindIt is a Python program for identifying a song given a \textbf{short noisy segment}, similar to Shazam
% \item An \textbf{audio fingerprint} consisting of constellations of major time-frequency peaks is stored in a hash table
% \item The music discovery program has \textbf{30+ stars} on GitHub
% \end{itemize}
% \end{itemize}
% \vspace{-10pt}
% %%%%%%%%%%%%%%%%%
% %%%%%%%%%%%%%%%%%
